\scanpage{7}
\mspage{5}
c) Si $\{u,v\}\in\mathcal P_2(\tilde I)$ tel que $\{u,v\}\cap\delta I=\varnothing$
(i.e. $\{u,v\}\in\mathcal P_2(\tilde I^\circ)$) alors conditions équivalentes

\unsure{1)} \deleted{\unsure{donc}} $\exists J\in S_I$ tel que $\delta J=\{u,v\}$

2) On a $v\in\tilde I_{x,u}$ \uline{ou} $v\in\tilde I_{y,u}$

3) On a $u\in\tilde I_{y,v}$ ou $u\in\tilde I_{x,v}$

\margin{\uline{NB} de \unsure{$=$|c}
\[
v\in\tilde I_{x,u}\Longleftrightarrow u\in\tilde I_{y,v}
\]
\[
v\in\tilde I_{y,u}\Longleftrightarrow u\in\tilde I_{x,v}
\]
\deleted{\unsure{Cela vaut pour} \unsure{$u,v\in\tilde I^\circ$}}
\deleted{\unsure{la condition de disjonction}}
\deleted{\unsure{$u\in\tilde I^\circ$, $v\in\tilde I^\circ$}}
\unsure{vaut} si on \deleted{\unsure{$u,v\in\tilde I$}} \unsure{suppose} que \deleted{\unsure{$u\neq v$}} $u,v\in\tilde I$ (quelconques, \unsure{supposons} $u,v\in\tilde I^\circ$ mais $u\neq v$)}

[\uline{NB} on dit alors que $u$ et $v$ sont \inserted{des lieux} \uline{disjoints} sur le segment $I$ -- on y reviendra\ldots]

\uline{Cor. 4} Si $\tilde I=\tilde J$ alors $I=J$ [donc on peut identifier $S$ à une partie de $\mathcal P(L)$]

$F_3)$ Soient $I,J\in S$, \inserted{$x\in L$} tels que $x\in\delta I\cap\delta J$.

Alors, \deleted{ou bien $\exists\unsure{z}\in I^\circ\cap J^\circ$ tel que $I_{x,z}=J_{x,z}$, ou bien} \deleted{$\unsure{\exists}$}
$\exists\unsure{u}\in I^\circ$, $v\in J^\circ$ tels que \deleted{$\unsure{I_{x,}}$} l'on ait
\[
\left\{
\begin{array}{l}
I_{x,u}=J_{x,v}\quad(\text{donc }u=v\text{ et }u=v\in I^\circ\cap J^\circ)\\
\uline{\text{ou}}\quad \tilde I_{x,u}\cap\tilde J_{x,v}=\{x\}
\end{array}
\right.
\]
\margin{\uline{Axiome des branches}}
\margin{Voir \unsure{renforcement} $F_3'$ \unsure{aussi} p.\ 10 et \deleted{\unsure{20}}}
\margin{écrire plutôt $u\in I\setminus\{x\}$, $v\in J\setminus\{x\}$ (si on n'admet pas condition de divisibilité) \deleted{\unsure{pour $I$}}}

\uline{NB} Ces deux relations sont mutuellement exclusives.

\deleted{La première, est}
La condition
\[
\exists z\in\tilde I^\circ\cap\unsure{\tilde J^\circ|J^\circ}
\quad\text{tel que}\quad
\unsure{I_{x,z}=I_{y,z}|I_{x,z}=J_{x,z}}
\]
est une relation \uline{d'équivalence} \inserted{$R_x$} dans l'ens \inserted{$S_x$} des segments $I\in S$ tels que $x\in\delta I$.
L'ensemble quotient
\[
B_x=\unsure{S|S_x}/R_x
\]
s'appelle l'\uline{ens des branches} en $x$, son cardinal l'ordre de $x$. Ordre 0:

\margin{\unsure{Ceci équivaut à} $\tilde I_{x,u}\cap\tilde I_{x,v}\neq\{x\}$.}
\margin{\unsure{Il manque} partie \textcircled{b} de $F_3$ ou \textcircled{$F_3'$}:
Soit $I\in S$, $\delta I=\{x,y\}$, $u,v\in I^\circ$, alors
$\exists w\in I^\circ$ tel que $u,v\in I_{x,w}^\circ$
(\unsure{i.e.|et} $\unsure{w\in I_{y,u}^\circ\cap I_{y,v}^\circ}$).
(donc $\exists w'\in I^\circ$ tel que $u,v\in I_{w',y}^\circ$).
\unsure{$w'\in I_{x,u}^\circ\cap I_{x,v}^\circ$}}
\mspage{6}

\rotatebox{90}{\unsure{EXTRA STRONG}}
\rotatebox{90}{\unsure{EXTRA STRONG}}
